
% \subsection{
%     جمع‌کننده‌های بیت نقلی آبشاری
%     \protect\efnote{Ripple Carry Adder (RCA)}
% }
% در این بخش با استفاده از بلوک‌های جمع‌کننده کامل بخش قبل (که ساختار مبتنی بر TG عملکرد بهتری ارائه داد و از آن استفاده می‌کنیم) به صورت سلسله‌مراتبی جمع‌کننده‌های 2بیتی، 4بیتی و درنهایت 8بیتی را با اتصال آبشاری بلوک‌های جمع‌کننده کامل طراحی کرد و به بررسی عملکرد و تحلیل آن‌ها می‌پردازیم.
% از آنجایی که بلوک‌ها به صورت سری به‌یکدیگر متصل شده‌اند و تاخیر بیشینه هر بلوک جمع‌کننده از خروجی بیت نقلی Cout آن ناشی می‌شود بیشترین تاخیر مدار هنگامی‌است که یک بیت نقلی از Cin مدار وارد و از Cout آن خارج بشود(به این معنی که تمام بلوک‌ها بیت نقلی خروجی تولید کنند). به صورت تخمینی انتظار داریم که تاخیر بیشینه مدار \lr{RCA-Nbit} تقریبا برابر باشد با:
% \lr{$$t_{RCA} \approx t_{A_0 \to Cout_0} + (N-2) \times t_{Cin \to Cout} + t_{Cin_7 \to Cout}$$}
% عملکرد مدار‌را با فایل \texttt{input.vec} بررسی می‌کنیم. همچنین برای بهتر به نمایش گذاشتن خروجی‌ها، بیت نقلی ورودی را صفر در نظر گرفته‌ایم.

% \vspace{1cm}

% \newimage[0.9\textwidth]{images/rca/RCA.jpg}{
%     شماتیک مدار جمع‌کننده بیت نقلی آبشاری 8 بیتی
% }{fig:rca}

% \newpage

% \newcode[0.9]{spice}{codes/rca/rca.sp}{مدل SPICE مدار جمع‌کننده بیت نقلی آبشاری}{code:rca}

% \newcode[0.9]{spice}{codes/rca/rca-vec.sp}{بخشی از فایل \protect\texttt{input.vec} مورد استفاده برای RCA 8بیت}{code:rca-vec}

% %%%%%%%%%%%%%%%%%%%%%%%%%%%%%%%%%%%%%%%%%%%%%%
% %%%%%%%%%%%%%%%%%%%%%%%%%%%%%%%%%%%%%%%%%%%%%%
% %%%%%%%%%%%%%%%%%%%%%%%%%%%%%%%%%%%%%%%%%%%%%%
% %%%%%%%%%%%%%%%%%%%%%%%%%%%%%%%%%%%%%%%%%%%%%%
% %%%%%%%%%%%%%%%%%%%%%%%%%%%%%%%%%%%%%%%%%%%%%%
% %%%%%%%%%%%%%%%%%%%%%%%%%%%%%%%%%%%%%%%%%%%%%%
% %%%%%%%%%%%%%%%%%%%%%%%%%%%%%%%%%%%%%%%%%%%%%%
% %%%%%%%%%%%%%%%%%%%%%%%%%%%%%%%%%%%%%%%%%%%%%%

% \newpage
% \subsubsection{عملکرد}
% بررسی عملکرد مدار 8بیتی ما با استفاده از \texttt{input.vec} صورت گرفته و تمام حالات ممکن برای ورودی‌ها به آن داده شده و شبیه‌سازی صورت گرفته است.
% همانطوری که از فایل \texttt{sim.lis} مشخص است پس از حدود 2 ساعت شبیه‌سازی و بدون هیچ خطایی در فایل \texttt{.errX} نشان از موفقیت بودن عملکرد مدار است.
% در ادامه بخشی از فایل \texttt{input.vec} و خروجی فایل \texttt{sim.lis} به اختصار آورده شده است.
% اما به دلیل خواسته‌ی پروژه اسکریپت پایتون پروژه اول را هم تغییر میدهیم تا بررسی توان و تحلیل مونت‌کارلو و و مابقی آزمایش‌ها را بر اساس آن انجام دهیم.
% \newcode[0.7]{spice}{codes/rca/rca-lis.sp}{خروجی فایل \protect\texttt{sim.lis} مربوط به شبیه‌سازی RCA 8بیت}{code:rca-lis}
% \vspace{-0.2cm}
% \newimage[0.9\textwidth]{images/rca/RCA-func-tree.png}{
%     همانطور که قابل مشاهده‌است حجم فایل \protect\texttt{input.vec} و همچنین \protect\texttt{sim.tr} بشدت بالا بوده از همین نظر آن‌ها در پروژه نگه داری نشدند.
% }{fig:rca-tree}


% %%%%%%%%%%%%%%%%%%%%%%%%%%%%%%%%%%%%%%%%%%%%%%
% %%%%%%%%%%%%%%%%%%%%%%%%%%%%%%%%%%%%%%%%%%%%%%
% %%%%%%%%%%%%%%%%%%%%%%%%%%%%%%%%%%%%%%%%%%%%%%
% %%%%%%%%%%%%%%%%%%%%%%%%%%%%%%%%%%%%%%%%%%%%%%
% %%%%%%%%%%%%%%%%%%%%%%%%%%%%%%%%%%%%%%%%%%%%%%
% %%%%%%%%%%%%%%%%%%%%%%%%%%%%%%%%%%%%%%%%%%%%%%
% %%%%%%%%%%%%%%%%%%%%%%%%%%%%%%%%%%%%%%%%%%%%%%
% %%%%%%%%%%%%%%%%%%%%%%%%%%%%%%%%%%%%%%%%%%%%%%
% %%%%%%%%%%%%%%%%%%%%%%%%%%%%%%%%%%%%%%%%%%%%%%
% %%%%%%%%%%%%%%%%%%%%%%%%%%%%%%%%%%%%%%%%%%%%%%
% %%%%%%%%%%%%%%%%%%%%%%%%%%%%%%%%%%%%%%%%%%%%%%
% %%%%%%%%%%%%%%%%%%%%%%%%%%%%%%%%%%%%%%%%%%%%%%
% %%%%%%%%%%%%%%%%%%%%%%%%%%%%%%%%%%%%%%%%%%%%%%
% %%%%%%%%%%%%%%%%%%%%%%%%%%%%%%%%%%%%%%%%%%%%%%
% %%%%%%%%%%%%%%%%%%%%%%%%%%%%%%%%%%%%%%%%%%%%%%

% \newpage

% \subsubsection{تاخیر زمانی}
% در ادامه تاخیر مدار به ازای تغییری\efnote{Transition} که پیش‌تر راجع‌به آن صحبت شد و منجر به بیشترین تاخیر در مدار ما می‌شود را در شبیه‌سازی سبک تری به نمایش می‌گذاریم.
% شبیه‌سازی های مربوط به این بخش در پوشه‌ی \texttt{delay} هستند.

% \newimage[0.78\textwidth]{images/rca/RCA-func.png}{
%     شبیه‌سازی RCA 8بیتی
% }{fig:rca-func}

% \newimage[0.78\textwidth]{images/rca/RCA-timing.png}{
%     شبیه‌سازی RCA 8بیتی با تمرکز بر بیشترین تاخیر در مدار که همانطور که قابل مشاهده هست به صورت آبشاری از انتقال بیت نقلی در مدار ایجاد می‌شود
% }{fig:rca-timing}

% تاخیر‌های اندازه‌گیری شده به شرح نمودار زیر است.
% \newimage[0.9\textwidth]{images/rca/RCA-delay.png}{
%     تاخیر‌های اندازه‌گیری شده برای \lr{RCA-8bit} در گوشه‌های مختلف طراحی در دمای 25 درجه
% }{fig:rca-delay}

% %%%%%%%%%%%%%%%%%%%%%%%%%%%%%%%%%%%%%%%%%%%%%%
% %%%%%%%%%%%%%%%%%%%%%%%%%%%%%%%%%%%%%%%%%%%%%%
% %%%%%%%%%%%%%%%%%%%%%%%%%%%%%%%%%%%%%%%%%%%%%%
% %%%%%%%%%%%%%%%%%%%%%%%%%%%%%%%%%%%%%%%%%%%%%%
% %%%%%%%%%%%%%%%%%%%%%%%%%%%%%%%%%%%%%%%%%%%%%%
% %%%%%%%%%%%%%%%%%%%%%%%%%%%%%%%%%%%%%%%%%%%%%%
% %%%%%%%%%%%%%%%%%%%%%%%%%%%%%%%%%%%%%%%%%%%%%%
% %%%%%%%%%%%%%%%%%%%%%%%%%%%%%%%%%%%%%%%%%%%%%%
% %%%%%%%%%%%%%%%%%%%%%%%%%%%%%%%%%%%%%%%%%%%%%%
% %%%%%%%%%%%%%%%%%%%%%%%%%%%%%%%%%%%%%%%%%%%%%%
% %%%%%%%%%%%%%%%%%%%%%%%%%%%%%%%%%%%%%%%%%%%%%%
% %%%%%%%%%%%%%%%%%%%%%%%%%%%%%%%%%%%%%%%%%%%%%%
% %%%%%%%%%%%%%%%%%%%%%%%%%%%%%%%%%%%%%%%%%%%%%%
% %%%%%%%%%%%%%%%%%%%%%%%%%%%%%%%%%%%%%%%%%%%%%%

% \newpage
% \subsubsection{توان مصرفی}
% برای بررسی توان متوسط مصرفی مدار از یک فایل \texttt{input.vec} به عنوان سمپلی رندوم استفاده کرده‌ایم.
% برای بررسی توان مصرفی متوسط از دستور \texttt{.MEASURE TRAN P\_avg AVG POWER} استفاده می‌کنیم.
% همچنین فایل تولید ورودی رندوم را به صورتی تغییر دادیم که برای محاسبه‌ی توان استاتیک آن بازه‌ی ثابت بودن خروجی‌ها را تنها درنظر بگیرد.
% به این ترتیب برداری از توان‌های متوسط اندازه‌گیری شده در هر اعمال ورودی به مدار خواهیم داشت که به دلیل ثابت بودن خروجی‌ها و معتبر بودن آنها تنها توان استاتیک مدار را شامل می‌شود.
% بردار ورودی‌ها 100 تست است.

% در شبیه‌سازی ها برای گوشه‌ی SS به دلیل افزایش تاخیر زمانی با این خطا مواجه‌شدیم:
% \newcode[]{spice}{codes/rca/rca-err-SS.sp}{خطای فایل \texttt{input.vec} برای گوشه‌ی طراحی SS بخش توان}{code:err-SS}

% \newimage[\textwidth]{images/rca/RCA-power.png}{
%     توان متوسط اندازه‌گیری شده برای \lr{RCA-8bit} در گوشه‌های مختلف طراحی در دمای 25 درجه
% }{fig:rca-power}

% به دلیل بزرگ درنظر گرفتن عرض ترانزیستور‌ها (\lr{$1\mu m$}) ترانزیستور‌های ما \lr{Long Channel} هستند.
% به این ترتیب بیشتر توان مصرفی مدار مربوط به \lr{Transition}ها خواهد بود نه جریان‌های نشتی.
% همانطور که در نمودار‌ مقایسه مصرف‌ توان نیز قابل مشاهده‌است سهم توان استاتیک در حدود نیم درصد است.
% همچنین قابل مشاهده‌است که هرچه به گوشه‌های سریع‌تر می‌رویم توان استاتیم به صورت نمایی افزایش می‌یابد که این به دلیل رابطه‌ی نمایی جریان نشتی با
% ولتاژ آستانه دارد و در گوشه‌های سریع‌تر این آستاه کمتر است:
% $$I_{sub} \propto e^{\frac{-V_{th}}{n v_T}}$$

% % برای اثبات اینکه با نحوه کارکرد فایل \texttt{.vec} آشنایی پیدا کرده‌ایم این خروجی دارای ارور را در گزارش قرار می‌دهیم:
% % \newcode[0.8]{spice}{codes/err.sp}{نمونه خروجی دارای خطای فایل \texttt{input.vec} که در بخش توان حاصل شد}{code:err}


% %%%%%%%%%%%%%%%%%%%%%%%%%%%%%%%%%%%%%%%%%%%%%%
% %%%%%%%%%%%%%%%%%%%%%%%%%%%%%%%%%%%%%%%%%%%%%%
% %%%%%%%%%%%%%%%%%%%%%%%%%%%%%%%%%%%%%%%%%%%%%%
% %%%%%%%%%%%%%%%%%%%%%%%%%%%%%%%%%%%%%%%%%%%%%%
% %%%%%%%%%%%%%%%%%%%%%%%%%%%%%%%%%%%%%%%%%%%%%%
% %%%%%%%%%%%%%%%%%%%%%%%%%%%%%%%%%%%%%%%%%%%%%%
% %%%%%%%%%%%%%%%%%%%%%%%%%%%%%%%%%%%%%%%%%%%%%%
% %%%%%%%%%%%%%%%%%%%%%%%%%%%%%%%%%%%%%%%%%%%%%%
% %%%%%%%%%%%%%%%%%%%%%%%%%%%%%%%%%%%%%%%%%%%%%%
% %%%%%%%%%%%%%%%%%%%%%%%%%%%%%%%%%%%%%%%%%%%%%%
% %%%%%%%%%%%%%%%%%%%%%%%%%%%%%%%%%%%%%%%%%%%%%%
% %%%%%%%%%%%%%%%%%%%%%%%%%%%%%%%%%%%%%%%%%%%%%%
% %%%%%%%%%%%%%%%%%%%%%%%%%%%%%%%%%%%%%%%%%%%%%%
% %%%%%%%%%%%%%%%%%%%%%%%%%%%%%%%%%%%%%%%%%%%%%%

% \subsubsection{تاثیرات دما بر پارامتر‌های مدار}

% \newimage[0.79\textwidth]{images/rca/RCA-power-peak-temp.png}{
%     توان پیک (بیشینه) اندازه‌گیری شده برای \lr{RCA-8bit} در گوشه‌ی TT برای دما‌های کاری مختلف
% }{fig:rca-power-peak-temp}

% پیک توان معمولاً در لحظاتی رخ می‌دهد که مدار بیشترین فعالیت سوئیچینگ را دارد و خازن‌ها در حال شارژ و دشارژ شدن هستند. این توان از نوع توان دینامیک سوئیچینگ و مشخص‌تر از نوع توان اتصال کوتاه است.

% ازانجایی که جریان مدار در حالت روشن با افزایش دما کاهش پیدا می‌کند در نتیجه آن توان لحظه‌ای پیک که حاصل ضرب $V_{dd} * I_{on}$ هست کاهش پیدا می‌کند.


% \newimage[0.79\textwidth]{images/rca/RCA-power-temp.png}{
%     توان متوسط اندازه‌گیری شده برای \lr{RCA-8bit} در گوشه‌ی TT برای دما‌های کاری مختلف
% }{fig:rca-power-temp}

% با افزایش دما، ولتاژ آستانه$V_{th}$ ترانزیستورها کاهش می‌یابد .
% این در کنار افزایش ولتاژ حرارتی ($v_T$) سبب می‌شود که جریان نشتی استاتیک مدار به صورت نمایی با افزایش دما رشد کند.
% از آنجایی که توان نشتی در تمام لحظات به طور مداوم مصرف می‌شود و اینکه توان پیک در لحاظ کمی بوجود می‌آید میانگین توان مصرفی کل را بالا می‌برد.

% \newimage[0.7\textwidth]{images/rca/RCA-power-time.png}{
%     توان لحظه‌ای اندازه‌گیری شده برای \lr{RCA-8bit} در گوشه‌ی TT در دمای 25 درجه
% }{fig:rca-power-time}


% \newimage[\textwidth]{images/rca/RCA-delay-temp.png}{
%     تاخیر‌های اندازه‌گیری شده برای \lr{RCA-8bit} در گوشه‌ی طراحی TT در دما‌های مختلف
% }{fig:rca-delay-temp}
% همانطور که قابل انتظار بود با افزایش دما تحرک‌پذیری حامل‌های جریان کاهش پیدا می‌کند که باعث افزایش تاخیر مدار می‌شود.
% هر چه زمان Transition سیگنال ها طولانی‌تر شود، مدت زمانی که هر دو ترانزیستور هم‌زمان روشن می‌مانند افزایش می‌یابد.
% این امر به طور غیرمستقیم باعث افزایش توان دینامیک اتصال کوتاه می‌شود.
% ازانجایی که در این مدار خروجی‌ها وابسته به خروجی‌های بخش‌های قبلی هستند این باعث اثر‌گذاری تاخیر بر توان مصرفی می‌شود.


%%%%%%%%%%%%%%%%%%%%%%%%%%%%%%%%%%%%%%%%%%%%%%
%%%%%%%%%%%%%%%%%%%%%%%%%%%%%%%%%%%%%%%%%%%%%%
%%%%%%%%%%%%%%%%%%%%%%%%%%%%%%%%%%%%%%%%%%%%%%
%%%%%%%%%%%%%%%%%%%%%%%%%%%%%%%%%%%%%%%%%%%%%%
%%%%%%%%%%%%%%%%%%%%%%%%%%%%%%%%%%%%%%%%%%%%%%
%%%%%%%%%%%%%%%%%%%%%%%%%%%%%%%%%%%%%%%%%%%%%%
%%%%%%%%%%%%%%%%%%%%%%%%%%%%%%%%%%%%%%%%%%%%%%
%%%%%%%%%%%%%%%%%%%%%%%%%%%%%%%%%%%%%%%%%%%%%%
%%%%%%%%%%%%%%%%%%%%%%%%%%%%%%%%%%%%%%%%%%%%%%
%%%%%%%%%%%%%%%%%%%%%%%%%%%%%%%%%%%%%%%%%%%%%%
%%%%%%%%%%%%%%%%%%%%%%%%%%%%%%%%%%%%%%%%%%%%%%
%%%%%%%%%%%%%%%%%%%%%%%%%%%%%%%%%%%%%%%%%%%%%%


\subsubsection{مقاومت در برابر نوسانات ولتاژ}
ولتاژ منبع تغذیه یکی از کلیدیترین پارامترهای تعیین کننده در میزان تأخیر و توان مصرفی مدار است.
برای ارزیابی میزان حساسیت مدار به نوسانات ولتاژ، از تحلیل آماری مونت‌کارلو (\lr{Monte Carlo}) استفاده می‌شود.
در این رویکرد، به پارامترهای مدار نه به عنوان مقادیر ثابت، بلکه به عنوان متغیرهای تصادفی با یک
تابع توزیع احتمال مشخص درنظر گرفته می‌شوند.
برای مدلسازی نوسانات ولتاژ تغذیه، ابتدا یک متغیر تصادفی با توزیع گاوسی (میانگین صفر و انحراف معیار 1.0)
به صورت زیر تعریف می‌کنیم.\\
سپس مقدار نهایی ولتاژمنبع تغذیه را (با فرض ولتاژ نامی 1 ولت برای تکنولوژی مورد نظر) به صورت ترکیبی
از مقدار نامی و مقدار نوسان تصادفی تعریف می‌کنیم.\\
در نهایت شبیه‌سازی مونت‌کارلو برای 500 تکرار و در قالب تحلیل گذرا با دستور زیر اجرا می‌شود:

\begin{LTR}
    \noindent
    \texttt{.param  dev=agauss(0,0.1)}\\
    \texttt{.param Vdd\_val=\textcolor{orange}{'1 + dev'}}\\
    \texttt{.tran 1p 20n SWEEP MONTE=500}
\end{LTR}

با اجرای این دستور، شبیه‌سازی 500 بار تکرار شده و در هر تکرار، نمونه‌ی جدیدی از متغیر تصادفی dev
تولید و اعمال می‌شود. پس از اتمام شبیه‌سازی، دو فایل خروجی کلیدی ایجاد خواهد شد: \\
فایل با پسوند \texttt{.mcX} حاوی مقادیر تولیدشده برای پارامتر dev در هر یک از 500 تکرار\\
فایل با پسوند \texttt{.mtX} حاوی مقادیر اندازه‌گیری شده به ازای هر 500 تکرار شبیه‌سازی

نکته مهم: مقادیر ثبت‌ شده در فایل \texttt{.mcX} برای توزیع گاوسی، به صورت نرمال شده نسبت به انحراف معیار
گزارش می‌شوند. بنابراین برای استخراج مقدار واقعی نوسان، باید اعداد این فایل را در مقدارانحراف معیار
ضرب کنیم.

درخت فایل‌های خروجی برای تحلیل \lr{Monte Carlo} به صورت زیر هست:
\newcode[1]{spice}{codes/rca/rca-monte-carlo-tree.sp}{درخت فایل‌های خروجی شبیه‌سازی مونت‌کارلو}{code:monte-carlo-tree}

